\section{Open Exposures}
\label{app:open}

This appendix lists the experiments and analyses that we did not
close before this draft.  None of the items below change the
headline numbers, but they are flagged for revision and for
reviewers who want to know what is missing.

\paragraph{No multi-seed variance on the headline numbers.}  All
evaluations are at fixed seed $0$.  Reviewer-2-style demand for
error bars across seeds is unmet.  Estimated cost: $\sim \$0.50$
for $3$ seeds $\times$ $4$ headline conditions ($N{=}200$ each).
At $N{=}200$, the exact binomial $95\%$ CI for an accuracy in the
$70\text{--}80\%$ range spans roughly $\pm 6\text{--}7$\,pp; we
report the per-condition CIs in every body table and in
Table~\ref{tab:prereg}.

\paragraph{ABL\_B sanity probe deferred.}  The $+0.0$\,pp result
for ABL\_B (Table~\ref{tab:distill_disentangle}; Track 1 v3 LoRA
$+$ commit-v3 with \code{COMMIT\_N\_BLOCKS=1}, full-response
training) is the most surprising single number in
\S\ref{subsec:distill_ablations}.  It is consistent with either
(i) the adapter firing as designed but having no effect on
extracted answers because the surgery lands downstream of the
already-pinned reasoning (the mechanism story we tell) or (ii) a
silent loader failure where the adapter does not modify inference
at all.  A small logits-diff probe ($5$ spot-check problems,
$\sim \$0.02$, $<10$\,min) would distinguish these.  We did not
run it before this draft; released adapter checkpoints permit
independent verification by readers.

\paragraph{Per-problem error-set diff for v1 vs.\ v2 not done.}
Track 2 v1 and v2 both produce $141/200$ correct on test
(\S\ref{subsec:distill_plateau}).  Whether the per-problem error
sets are identical or merely aggregate-equal is itself
informative: identical sets support a stronger ``the late-block
bottleneck binds the failure mode to specific problems''
reading, while non-identical sets support a weaker ``different
adapters fail and pass on different problems but the late-block
bottleneck caps both at the same total'' reading.  The body
makes the weaker claim.

\paragraph{LCH-centroid mechanism figure for Track 1 was not
produced.}  A latent-cluster-heatmap-style figure separating
prefix-robust LoRA representations from base \LLaDA{}
representations would harden the structural-separation thesis
from behavioural to mechanistic.  The JVP feasibility spike
failed because PEFT~\citep{mangrulkar2022peft} wraps LoRA in a
\code{lora\_magnitude\_vector} ModuleDict that requires a custom
JVP shim.  Deferred to a follow-up.

\paragraph{Track 1 v1 catastrophe is documented in
\code{e2/RESULTS\_TRACK1.md} but not in the body.}  v1 used $r{=}16$,
$\alpha{=}32$, learning rate $1\mathrm{e}{-}4$, $p_\text{mask}
\sim U(\epsilon, 1{-}\epsilon)$, and $14$k steps; the resulting
adapter regressed \ctwo{} from $74\%$ to $59.5\%$ ($-14.5$\,pp)
via bigram-repetition mode collapse (the model learned a
copy-back-context heuristic at training time and defaulted to it
at inference).  v2 fixed this with four hyperparameter changes
(smaller $r$, gentler LR, narrower $p_\text{mask}$ range, fewer
steps).  We mention v1 only in passing in
\appref{app:hyperparameters}; the contrast with v2 is a
publishable masked-diffusion SFT pitfall on its own and is
preserved in the source draft for a follow-up.

\paragraph{No cross-task transfer.}  All experiments are on
GSM8K-test~\citep{cobbe2021gsm8k}.  The diversity-expansion
finding in \S\ref{sec:axis3} in particular is suspicious-looking
on a single benchmark and deserves a multi-task replication
(MATH-Easy, ARC-Challenge are the natural next targets).

\paragraph{No comparison to BD3-LMs and DiffuLLaMA on the same
eval.}  The literature claims we discuss in
\S\ref{sec:discussion} are reconciled \emph{by reading} other
authors' results through the three-axis frame, not by re-running
their adapters under our prefix conditions.  A formal multi-DDLM
ablation (BD3-LMs~\citep{arriola2025bd3lms},
DiffuLLaMA~\citep{gong2024diffullama}, \LLaDA{}, possibly hybrid
multimodal models~\citep{zhou2024transfusion,xie2024showo,ma2025janusflow})
on the same prefix-tier dataset would be the cleanest
falsification target for the three-axis framing.  Out of scope
for this draft.

\paragraph{Diversity-mechanism test deferred.}  Whether the
diversity uptick in \S\ref{sec:axis3} is an
implicit-temperature effect or a real content-diversity effect
is open.  The cheap test (base \LLaDA{} at $t \in \{0.8, 0.9,
1.0\}$ on the same eval, read the $5/5$-same rate) was deferred
to revision.

\paragraph{Reranker baseline not run.}  A small classifier over
$5$ \cmaj{} branches as a compute-time alternative to v3's
param-time commit would test whether weight-space distillation
is uniquely useful.  Several days of work; deferred to a
follow-up paper.
